\addcontentsline{toc}{subsubsection}{Solution 60. The Squared-Logarithm Trap}

\begin{dlsummary}[title={\bfseries Solution 60. The Squared-Logarithm Trap},breakable]

\begin{proof}
The unique solution is
\[
\boxed{
f(x)=2^{(\log_2 x)^2}
\qquad(x>0).
}
\]

\paragraph{1. Translate multiplication into addition.}

Define
\[
g(t)=\log_2 f(2^t)\qquad(t\in\mathbb R).
\]
The range and normalisations give
\[
g(t)\ge0,\qquad g(1)=1,\qquad
g(t)\le t\quad(0\le t\le1).
\tag{3}
\]

The two functional inequalities become
\[
g(u+v)+g(u-v)\le2g(u)+2g(v), \tag{4}
\]
and
\[
g(2u)\ge4g(u). \tag{5}
\]

\paragraph{2. Recover evenness and exact scaling.}

At \(u=0\), (5) and nonnegativity force \(g(0)=0\).

Taking \(u=0\) in (4),
\[
g(v)+g(-v)\le2g(v),
\]
so \(g(-v)\le g(v)\). Replacing \(v\) by \(-v\) proves
\[
g(-v)=g(v).
\]

Taking \(v=u\) in (4), then combining with (5), gives
\[
\boxed{g(2u)=4g(u).} \tag{6}
\]

\paragraph{3. Turn the inequality into the parallelogram equation.}

Apply (4) to \(u+v,u-v\):
\[
g(2u)+g(2v)\le2g(u+v)+2g(u-v).
\]
Using (6), this reverses (4). Hence
\[
\boxed{
g(u+v)+g(u-v)=2g(u)+2g(v).
}
\tag{7}
\]

\paragraph{4. Determine rational values.}

For fixed \(t\), equation (7) gives
\[
g((n+1)t)+g((n-1)t)=2g(nt)+2g(t).
\]
Induction from \(g(0)=0\) proves
\[
g(nt)=n^2g(t)
\]
for every integer \(n\).

Since \(g(1)=1\), it follows that
\[
\boxed{g(q)=q^2\qquad(q\in\mathbb Q).} \tag{8}
\]

\paragraph{5. Eliminate discontinuous possibilities.}

From evenness and (3),
\[
0\le g(h)\le |h|\qquad(|h|\le1).
\]
Thus \(g(h)\to0\) as \(h\to0\).

We now transfer this control to any fixed input \(x\). For any positive integer \(N\), use (7) with inputs \(x/N\) and \(h\). Nonnegativity and integer scaling give
\[
g(x+Nh)\le2g(x)+2N^2g(h).
\]

Also, (6) and (7) imply midpoint convexity:
\[
g\!\left(\frac{a+b}{2}\right)
\le\frac{g(a)+g(b)}2.
\]
For \(N\) a power of two, repeated midpoint convexity therefore gives
\[
\begin{aligned}
g(x+h)
&\le\left(1-\frac1N\right)g(x)
+\frac1N g(x+Nh)\\
&\le g(x)+\frac{g(x)}N+2Ng(h).
\end{aligned}
\]
The same upper bound holds for \(g(x-h)\). Since (7) gives
\[
g(x+h)+g(x-h)\ge2g(x),
\]
we conclude
\[
|g(x+h)-g(x)|
\le\frac{g(x)}N+2Ng(h).
\]
Choose \(N\) large first, then \(h\) small. This proves continuity at every \(x\).

Now (8) and rational approximation force
\[
g(t)=t^2\qquad(t\in\mathbb R).
\]
Undoing the transformation yields the claimed formula.

\paragraph{6. Verify.}

Write \(a=\log_2 x\) and \(b=\log_2 y\). Then
\[
(a+b)^2+(a-b)^2=2a^2+2b^2,
\]
so (1) holds with equality. Likewise, (2) holds with equality.

Finally, if \(1\le t\le2\), then \(s=\log_2t\in[0,1]\), so \(s^2\le s\). Therefore
\[
f(t)=2^{s^2}\le2^s=t.
\]
All conditions are satisfied.
\end{proof}

\paragraph{Takeaways.}
\begin{itemize}
\item **Transform both inputs and outputs:** Logarithms expose the additive structure behind a multiplicative inequality.
\item **Change coordinates to reverse a bound:** The substitution \((u+v,u-v)\) forces equality.
\item **Rational values are not the finish line:** A separate regularity argument is essential.
\item **A local condition can have global force:** The bound on \([1,2]\) eliminates all irregular quadratic solutions.
\end{itemize}

\end{dlsummary}
